\section{Results}
\label{sec:results}

The evaluation covers forecast accuracy and the executed schedules that produce
it. Mean forecast loss and macro-averaged normalized forecast loss are co-primary
endpoints. The validation-selected fixed-schedule baseline is the structural
primary comparator, and the warning-rule baseline is the practical adaptive
comparator. Scalar-objective and algorithm comparisons concern complete training
configurations rather than isolated components.

\subsection{Forecasting performance on event-balanced evaluation windows}
\label{sec:main_result}

\Cref{tab:clean_main_comparisons} reports the primary paired analysis for the 22
post-selection evaluation seeds (119--140). The complete PD-PPO training
configuration has lower mean forecast loss and lower
macro-averaged normalized forecast loss than the validation-selected
fixed-schedule baseline in all 22 seeds. Mean forecast
loss decreases from 1.5744 to 1.4073, a relative reduction of 10.6\%. The
macro-averaged loss decreases from 1.0440 to 0.9602, a relative reduction of
8.0\%. The mean paired differences against the fixed-schedule baseline are +0.1671 for
mean forecast loss and +0.0838 for the macro-averaged loss, with 95\%
percentile bootstrap confidence intervals for the mean paired differences that
remain above zero. The exact one-sided sign test probability for a direction
fixed as lower loss, conditional on 22 independent signs under a 0.5 null, is
$2.38\times10^{-7}$ for either co-primary endpoint.

Seeds 117 and 118 supplied selected test-window metrics for a two-variant pilot
architecture comparison, after which the selected configuration was frozen.
They are excluded from the primary analysis. Supplementary Table S12 retains the
24-seed descriptive aggregate, including the 10.1\% and 7.9\% reductions, without
assigning it post-selection inferential status.
The PD-PPO scheduler also has lower values of both losses than the AoI-priority,
round-robin, and random-feasible baselines in all 22 seeds.

\begin{table*}[htbp]
  \centering
  \caption{Baseline comparisons with the complete PD-PPO training configuration on the event-balanced evaluation windows over 22 post-selection evaluation seeds. A positive paired difference means that the baseline has higher loss than PD-PPO.}
  \label{tab:clean_main_comparisons}
  \footnotesize
  \setlength{\tabcolsep}{2.4pt}
  \begin{tabular}{@{}>{\raggedright\arraybackslash}p{0.22\textwidth}>{\raggedright\arraybackslash}p{0.36\textwidth}>{\raggedright\arraybackslash}p{0.36\textwidth}@{}}
    \toprule
    Baseline & \shortstack[l]{Mean forecast loss\\Seeds with lower loss;\\mean paired difference\\{[95\% bootstrap CI]}} & \shortstack[l]{Macro-averaged loss\\Seeds with lower loss;\\mean paired difference\\{[95\% bootstrap CI]}} \\
    \midrule
    Fixed-schedule baseline & 22 of 22; +0.1671 [+0.1231, +0.2164] & 22 of 22; +0.0838 [+0.0708, +0.0967] \\
    AoI-priority baseline & 22 of 22; +0.2619 [+0.2331, +0.2928] & 22 of 22; +0.1655 [+0.1430, +0.1898] \\
    Round-robin baseline & 22 of 22; +0.2571 [+0.2249, +0.2900] & 22 of 22; +0.1664 [+0.1427, +0.1917] \\
    Random-feasible baseline & 22 of 22; +0.2681 [+0.2366, +0.3019] & 22 of 22; +0.1684 [+0.1446, +0.1933] \\
    Privileged one-step look-ahead reference using next-step test targets & 22 of 22; +0.2746 [+0.2391, +0.3107] & 22 of 22; +0.1855 [+0.1559, +0.2173] \\
    Warning-rule baseline & 8 of 22; -0.0039 [-0.0102, +0.0019] & 9 of 22; +0.0011 [-0.0049, +0.0081] \\
    Privileged true-label baseline & 9 of 22; -0.0042 [-0.0108, +0.0011] & 10 of 22; +0.0017 [-0.0045, +0.0086] \\
    \bottomrule
  \end{tabular}
  \begin{minipage}{0.98\textwidth}
  \vspace{0.4ex}\footnotesize Each interval is a 95\% bootstrap confidence interval for the paired difference. The validation-selected fixed-schedule baseline is the structural primary comparator; the warning-rule baseline is the practical adaptive comparator. The one-step look-ahead reference uses next-step test targets, and the true-label baseline uses the exact sequence $c_t,\ldots,c_{t+16}$; both receive information unavailable during deployment, and neither is an upper bound. The warning-rule baseline uses the same simulator-generated warning signals available to PD-PPO and a mapping selected on the calibration/validation partition.
  \end{minipage}
\end{table*}


The eight separately reset 512-step windows are constructed from test truth to
balance particle, flux, and thermal coverage and emphasize transport intensity.
The resulting event-balanced evaluation windows do not represent the natural
test distribution.
Supplementary Table S7 tests whether that
sampling rule drives the result by evaluating the same final policy checkpoints
and the fixed channel subsets selected on the calibration/validation partition once
over all 5,242 scoreable test time steps in each continuous test partition. The
PD-PPO scheduler again has lower mean forecast loss and lower macro-averaged loss
in all 22 seeds. The mean paired differences are +0.1347 for mean forecast loss and +0.0836 for the
macro-averaged loss, and both 95\% percentile bootstrap confidence intervals for
the mean paired differences remain above zero. The
observed paired advantage therefore extends beyond the windows emphasizing transport periods
used for the primary balanced comparison.
\Cref{fig:clean_main_evidence} summarizes the seed-level paired differences for
the main comparison with the fixed-schedule baseline and the additional baselines.

\begin{figure}[htbp]
  \centering
  \includegraphics[width=\linewidth]{figure_clean_main_evidence.pdf}
  \caption{Seed-level paired differences over 22 post-selection evaluation seeds. Panel A
  compares PD-PPO with the fixed-schedule baseline. Panel B compares the
  AoI-priority, round-robin, and random-feasible baselines and the privileged
  one-step look-ahead reference using next-step test targets.
  Panel C compares the warning-rule and privileged true-label baselines and the Double DQN
  training configuration.
  Positive values mean that the baseline has higher loss than PD-PPO.}
  \label{fig:clean_main_evidence}
\end{figure}

\subsection{Channel activation across operating conditions}
\label{sec:regime_allocation}

The aggregate gain is distributed across all three forecast event types
(\Cref{tab:clean_regime_decomposition}). Mean paired differences are positive
for particle, flux, and thermal event types, and each 95\% percentile bootstrap
confidence interval for the mean paired difference remains above zero. The
PD-PPO scheduler has lower loss in 18 of 22 particle comparisons, 12 of 22
flux comparisons, and 15 of 22 thermal comparisons. Equal weighting combines
the three distinct specialist-channel requirements and yields lower
macro-averaged loss in all 22 seeds.

\begin{table}[htbp]
  \centering
  \caption{Event-type-specific forecast losses and the macro-averaged normalized forecast loss. Lower values are better; a positive paired difference favors PD-PPO over the fixed-schedule baseline.}
  \label{tab:clean_regime_decomposition}
  \scriptsize
  \setlength{\tabcolsep}{2.6pt}
  \begin{tabular}{@{}lcccc@{}}
    \toprule
    Event type & PD-PPO & Fixed baseline & \shortstack{Seeds with\\lower loss} & \shortstack{Mean paired difference\\{[95\% bootstrap CI]}} \\
    \midrule
    Particle event type & 1.0749 & 1.1037 & 18 of 22 & +0.0287 [+0.0148, +0.0473] \\
    Flux event type & 0.7639 & 0.8930 & 12 of 22 & +0.1291 [+0.0642, +0.1938] \\
    Thermal event type & 1.0418 & 1.1353 & 15 of 22 & +0.0935 [+0.0444, +0.1452] \\
    Macro-averaged loss & 0.9602 & 1.0440 & 22 of 22 & +0.0838 [+0.0708, +0.0967] \\
    \bottomrule
  \end{tabular}
\end{table}


\Cref{fig:clean_mechanism} compares optional-channel activation by PD-PPO and the
validation-selected fixed-schedule baseline. The laser disdrometer channel is active for
99.2\% of particle-event steps, the FC4 snow mass-flux channel for 97.8\% of
flux-event steps, and the infrared snow-surface-temperature channel for 99.7\% of
thermal-event steps. Calm periods use a mixture of the four optional channels with
nonzero test activation frequency. Calibration/validation selects the FC4 snow
mass-flux channel in 11 seeds, the infrared snow-surface-temperature channel in
nine, and the laser disdrometer channel in two, after which the selected subset
remains fixed on the test partition.

\begin{figure}[H]
  \centering
  \includegraphics[width=\linewidth]{figure_clean_mechanism.pdf}
  \caption{Optional-channel activation by operating condition. Panel A shows the
  shielded temperature--humidity, infrared snow-surface-temperature, laser
  disdrometer, and FC4 snow mass-flux channels. Panel B shows the fixed channels
  selected on the calibration/validation partition. Panel C reports executed-action
  entropy, condition--executed-action mutual information, and mean normalized
  Hamming distance between adjacent executed masks.}
  \label{fig:clean_mechanism}
\end{figure}

All 22 executed-action sequences pass the behavior-complexity gate; none is
fixed-like or simple-cycle-like under its prespecified criteria. Mean
executed-action entropy is 1.872 bits, and mean condition--executed-action mutual
information is 0.632 bits. The mean normalized Hamming distance between adjacent
executed masks is 0.00369. The meteorological core channel remains active by
construction, while the executed schedules do not activate the basic radiometer
channel in the selected test windows. At least three of the four active optional
channels have intermediate channel activation frequency in every seed; all four
do so in 21 of 22 seeds.
No channel was deactivated before completing warm-up because $D_{\min}=6$
exceeds the maximum two-step warm-up. Executed channel changes occur mainly when
the operating condition changes.
The gate's fixed-like and simple-cycle-like criteria exclude fixed schedules and
simple cycles. These behavior diagnostics concatenate the
eight separately reset windows and include the seven reset boundaries; they are
descriptive diagnostics, not a primary performance endpoint.

\subsection{Scalar-objective and training-configuration comparisons}
\label{sec:objective_learner_controls}

The AoI-objective and uncertainty-objective PPO configurations retain the PPO
architecture, scheduler input, feasible-action mask, supervised target actions
used only during training, auxiliary future-condition classifier, chronological
partitions, and frozen primary forecaster, but use AoI or a bounded diagonal
uncertainty proxy as the scalar rollout objective. Because each objective supplies its own advantages,
the comparison also changes the PPO surrogate, value target, and
advantage-weighted behavior cloning weights. \Cref{tab:clean_learning_controls}
reports the paired differences from PD-PPO. The mean macro-averaged loss
difference is +0.0015 for the AoI configuration, with lower loss for the
forecast-loss-objective configuration in 12 of 22 seeds, and +0.0013 for the
uncertainty configuration, with lower loss in 12 of 22 seeds. For both
alternative configurations, both 95\% bootstrap confidence intervals for the
paired difference in macro-averaged loss and mean forecast loss include zero.
These comparisons establish neither a population-level directional difference
nor equivalence among scalar objectives. All three matched PPO configurations
nevertheless have lower macro-averaged loss than the fixed-schedule baseline in
all 22 seeds and meet the channel-use acceptance criterion in every seed.
Warm-up interruption is structurally zero for all
three configurations because $D_{\min}=6$ exceeds the maximum two-step warm-up.

The Double DQN training configuration uses the same scheduler observation,
reward based on forecast loss, action set, constraints, policy-training partition, and frozen primary
forecaster. PD-PPO has lower
macro-averaged loss in all 22 seeds and lower mean forecast loss in 21 of 22,
with mean baseline-minus-PD-PPO paired differences of +0.0706 and +0.1423,
respectively. Double DQN has lower macro-averaged loss than the fixed-schedule baseline
in 12 of 22 comparisons and meets the channel-use acceptance criterion
in 5 of 22 seeds; no channel can be deactivated before completing warm-up under
the reported $D_{\min}$. All test metrics use the same frozen primary forecaster.
This is a training-configuration comparison on the same scheduling problem,
because the complete PD-PPO configuration uses behavior cloning pretraining,
the continuing advantage-weighted behavior cloning term, and the auxiliary
future-condition classifier; it is not an isolated comparison of
the PPO and Double DQN update rules.

\begin{table*}[htbp]
  \centering
  \caption{Scalar-objective PPO and Double DQN training-configuration comparisons. A positive paired difference indicates lower loss for PD-PPO.}
  \label{tab:clean_learning_controls}
  \scriptsize
  \setlength{\tabcolsep}{2.1pt}
  \begin{tabular}{@{}>{\raggedright\arraybackslash}p{0.14\textwidth}>{\raggedright\arraybackslash}p{0.27\textwidth}>{\raggedright\arraybackslash}p{0.27\textwidth}>{\centering\arraybackslash}p{0.12\textwidth}>{\centering\arraybackslash}p{0.14\textwidth}@{}}
    \toprule
    Comparison & Mean forecast loss\newline Seeds with lower loss;\newline mean paired difference\newline [95\% bootstrap CI] & Macro-averaged loss\newline Seeds with lower loss;\newline mean paired difference\newline [95\% bootstrap CI] & Lower macro-averaged\newline loss than fixed baseline & Channel-use acceptance\newline criterion \\
    \midrule
    AoI-objective PPO configuration & 10 of 22; -0.0003 [-0.0073, +0.0063] & 12 of 22; +0.0015 [-0.0037, +0.0066] & 22 of 22 & 22 of 22 \\
    Uncertainty-objective PPO configuration & 10 of 22; +0.0008 [-0.0063, +0.0076] & 12 of 22; +0.0013 [-0.0035, +0.0054] & 22 of 22 & 22 of 22 \\
    Double DQN training configuration & 21 of 22; +0.1423 [+0.1036, +0.1827] & 22 of 22; +0.0706 [+0.0535, +0.0875] & 12 of 22 & 5 of 22 \\
    \bottomrule
  \end{tabular}
  \begin{minipage}{0.98\textwidth}
  \vspace{0.4ex}\footnotesize Each interval is a 95\% bootstrap confidence interval. The objective rows compare matched PPO training configurations with different scalar objectives. Each objective changes the PPO surrogate, value target, advantage, and the weights in the advantage-weighted behavior cloning term; the rows are therefore not isolated scalar-reward ablations. The configurations otherwise share the scheduler input, training-only supervised targets, auxiliary future-condition classifier, candidate action set, and constraints. The Double DQN row is a training-configuration comparison: it shares the state representation, reward based on forecast loss, action set, and constraints, but omits behavior cloning pretraining, the continuing behavior cloning term, and the auxiliary classifier. All configurations have zero warm-up interruption structurally because $D_{\min}=6$ exceeds the maximum two-step warm-up; the final column is therefore determined by the channel-use acceptance criterion. All test metrics use the same frozen primary forecaster.
  \end{minipage}
\end{table*}


\subsection{Baselines using additional information}
\label{sec:information_references}

The privileged one-step look-ahead reference using next-step test targets uses
information unavailable during deployment. At each time step it reads next-step targets on
the test partition and evaluates feasible actions with their next-step forecast
loss. Its one-step scoring does not account for the longer-term consequences of
the minimum dwell-time constraint, warm-up, and observation age. The PD-PPO scheduler
has lower mean forecast loss and lower macro-averaged normalized forecast loss
in all 22 seeds
(\Cref{tab:clean_main_comparisons}). For this reference, the mean normalized
Hamming distance between adjacent executed masks is 0.0280. Despite
access to next-step forecast loss, it has higher mean forecast loss and higher
macro-averaged normalized forecast loss in every seed, and its executed activation indicators change approximately 7.6 times as often as those of
PD-PPO.

The warning-rule baseline maps simulated warning signals to proposed actions.
The privileged true-label baseline uses the exact sequence $c_t,\ldots,c_{t+16}$
and a condition-to-action mapping selected on the calibration/validation
partition. These condition labels are unavailable during deployment.
Their mean macro-averaged loss differences from PD-PPO are
+0.0011 and +0.0017, respectively. For both baselines, the 95\% percentile
bootstrap confidence intervals for the mean paired differences include zero
for macro-averaged loss and mean forecast loss. These intervals establish neither
a population-level directional difference nor equivalence. The true-label baseline is neither
deployable nor an upper bound on feasible sequential control.

\subsection{Alternative-forecaster analysis}
\label{sec:secondary_forecaster}

The alternative-forecaster analysis evaluates the test observation sequences
with a ridge forecaster. Fixed-schedule selection was repeated on the
calibration/validation partition before test evaluation. PD-PPO has lower
macro-averaged loss than the primary-forecaster fixed-schedule baseline in all 22
comparisons and than the ridge-forecaster fixed-schedule baseline in 21 of 22
(Supplementary Table S8).
The mean paired difference against the fixed
schedule selected using the ridge forecaster is +0.1339,
with a 95\% percentile bootstrap confidence interval for the mean paired
difference of [+0.1103, +0.1563]. This evaluation changes both
the forecaster architecture and the validation-selected fixed-schedule baseline without retraining the
policy.
\Cref{fig:clean_controls} summarizes the scalar-objective comparisons and the two
ridge-forecaster comparisons.

\begin{figure}[H]
  \centering
  \includegraphics[width=\linewidth]{figure_clean_controls.pdf}
  \caption{Scalar-objective training-configuration comparisons and
  alternative-forecaster analysis over 22 post-selection seeds. Panel
  A compares matched PPO configurations with different scalar objectives. Panel B gives macro-averaged loss differences for the
  primary-forecaster fixed-schedule baseline and the fixed schedule selected with
  the ridge forecaster. Positive values mean that the baseline has higher loss than
  PD-PPO.}
  \label{fig:clean_controls}
\end{figure}
